\section{Method}

\subsection{Problem Formulation}
Let $x \in \mathbb{C}^{N \times N \times N}$ denote the object-domain complex
volume. The measured diffraction modulus is
\begin{equation}
  y = \left|\mathcal{F}\{x\}\right|,
\end{equation}
where $\mathcal{F}$ denotes the 3D Fourier transform. The network receives $y$
and predicts object-domain amplitude $\hat{a}$ and phase $\hat{\phi}$.

\subsection{Network Output and Post-processing}
The predicted complex object is constructed as
\begin{equation}
  \hat{x} = s(\hat{a}) \cdot \hat{a}\exp(i\hat{\phi}),
\end{equation}
where $s(\hat{a})$ is the support mask. In the AutoPhaseNN-style pipeline, a
hard support threshold is applied:
\begin{equation}
  s(\hat{a}) =
  \begin{cases}
  1, & \hat{a} \geq T,\\
  0, & \hat{a} < T.
  \end{cases}
\end{equation}
The predicted diffraction modulus is
\begin{equation}
  \hat{y} = \left|\mathcal{F}\{\operatorname{ifftshift}(\hat{x})\}\right|.
\end{equation}
This layer is central to the comparison because it couples three quantities:
the predicted amplitude determines the support mask, the support mask gates the
complex object, and the gated object determines the far-field loss. A low
far-field error therefore does not directly imply that $\hat{a}$, $\hat{\phi}$,
or $s(\hat{a})$ are individually correct.

\subsection{Training Objective}
The main training objective follows the paper-style far-field modulus L1 loss:
\begin{equation}
  \mathcal{L}_{FT} = \frac{1}{|\Omega|}\sum_{v \in \Omega}
  \left| \hat{y}_v - y_v \right|.
\end{equation}
Additional real-space metrics are reported for diagnosis, but the default
AutoPhaseNN reproduction uses only the diffraction-domain loss.

\subsection{Controlled Backbone Replacement}
The UMamba experiment is designed as a controlled replacement rather than a new
pipeline. The intended controlled variables are:
\begin{itemize}
  \item identical train/validation sample selection;
  \item identical preprocessing of diffraction, amplitude, and phase tensors;
  \item identical physical post-processing from amplitude and phase to far-field
  modulus;
  \item identical far-field training loss unless an ablation explicitly changes
  it;
  \item identical evaluation metrics and visualization protocol.
\end{itemize}
Under this setup, any persistent difference in reconstruction quality is more
likely to come from the model architecture or from the strength of the encoded
spatial prior.

\subsection{Model Variants}
We compare the following models:
\begin{itemize}
  \item AutoPhaseNN reproduction with the original convolutional architecture.
  \item UMamba replacement with the same data and post-processing pipeline.
  \item UMamba variants with modified support thresholding or explicit spatial
  priors.
\end{itemize}
